%%%%%%%%%%%%%%%%%%%%%%%%%%%%%%%%%%%%%%%%%%%%%%%%%%%%%%%%%%%%%%%%%%%%%%%%%%%%%%%%
% Medium Length Graduate Curriculum Vitae
% LaTeX Template
% Version 1.2 (3/28/15)
%
% This template has been downloaded from:
% http://www.LaTeXTemplates.com
%
% Original author:
% Rensselaer Polytechnic Institute 
% (http://www.rpi.edu/dept/arc/training/latex/resumes/)
%
% Modified by:
% Daniel L Marks <xleafr@gmail.com> 3/28/2015
%
% Important note:
% This template requires the res.cls file to be in the same directory as the
% .tex file. The res.cls file provides the resume style used for structuring the
% document.
%
%%%%%%%%%%%%%%%%%%%%%%%%%%%%%%%%%%%%%%%%%%%%%%%%%%%%%%%%%%%%%%%%%%%%%%%%%%%%%%%%

%-------------------------------------------------------------------------------
%	PACKAGES AND OTHER DOCUMENT CONFIGURATIONS
%-------------------------------------------------------------------------------

%%%%%%%%%%%%%%%%%%%%%%%%%%%%%%%%%%%%%%%%%%%%%%%%%%%%%%%%%%%%%%%%%%%%%%%%%%%%%%%%
% You can have multiple style options the legal options ones are:
%
%   centered:	the name and address are centered at the top of the page 
%				(default)
%
%   line:		the name is the left with a horizontal line then the address to
%				the right
%
%   overlapped:	the section titles overlap the body text (default)
%
%   margin:		the section titles are to the left of the body text
%		
%   11pt:		use 11 point fonts instead of 10 point fonts
%
%   12pt:		use 12 point fonts instead of 10 point fonts
%
%%%%%%%%%%%%%%%%%%%%%%%%%%%%%%%%%%%%%%%%%%%%%%%%%%%%%%%%%%%%%%%%%%%%%%%%%%%%%%%%

\documentclass[margin]{res}  

% Default font is the helvetica postscript font
\usepackage{helvet}
\usepackage{graphicx}
\usepackage{hyperref}
\hypersetup{hidelinks}

\usepackage[utf8x]{inputenc}

% Increase text height
\textheight=700pt

\newenvironment{absolutelynopagebreak}
  {\par\nobreak\vfil\penalty0\vfilneg
   \vtop\bgroup}
  {\par\xdef\tpd{\the\prevdepth}\egroup
   \prevdepth=\tpd}



\begin{document}

%-------------------------------------------------------------------------------
%	NAME AND ADDRESS SECTION
%-------------------------------------------------------------------------------
\name{Konstantinos Kallas}

% Note that addresses can be used for other contact information:
% -phone numbers
% -email addresses
% -linked-in profile

% \address{Mpoumpoulinas 27 \\Dionisos, Greece, 14576\\\\Address 1 line 3}
% \address{Address 2 line 1\\Address 2 line 2\\Address 2 line 3}

% Uncomment to add a third address
%\address{Address 3 line 1\\Address 3 line 2\\Address 3 line 3}
%-------------------------------------------------------------------------------

\begin{resume}

%-------------------------------------------------------------------------------
%	EDUCATION SECTION
%-------------------------------------------------------------------------------
\section{Contact} 
%% Address: Mpoumpoulinas 27, Dionisos, Greece, 14576 \\
%% Phone: +30 6957814666 \\
Email: \href{mailto:kallas@seas.upenn.edu}{kallas@seas.upenn.edu} -- Website: \href{https://angelhof.github.io/}{angelhof.github.io}
%% Github: \href{https://github.com/angelhof}{angelhof} \\


\section{Education}
\textbf{University of Pennsylvania} \hfill {\em September 2018 -- present}\\
Computer and Information Science, PhD student \\
Advisor: Prof. Rajeev Alur

\textbf{Eighth Summer School on Formal Techniques} \hfill \mbox{\em May 2018}\\
Menlo College, Atherton, CA \\
Organised by SRI International

\textbf{National Technical University of Athens} \hfill \mbox{\em October 2012 -- February 2018}\\
Electrical and Computer Engineering, Diploma \\ %% \hfill {\em GPA: 8.37/10} \\
% Specialization: Computer Science \hfill {\em Major (Last 2 years) GPA: 9.05/10} \\
Thesis: \textit{``HiPErJiT: A Profile-Driven Just-in-Time Compiler for Erlang''} \\
Advisor: Prof. Kostis Sagonas 

\textbf{Universita degli Studi dell'Aquila} \hfill {\em February 2015 -- July 2015} \\
Computer Science, Erasmus Mobility Program %% \hfill {\em GPA: 9.4/10} \\
% Selected Courses: Programming Languages, Cryptography, and Software Engineering

% \textbf{Dionisos 1st High School, Greece} \hfill {\em Graduation: July 2012} \\
% Greek Apolitirion, ranked first among 120 students \hfill {\em GPA: 19.4/20}



\section{Working Experience}

\textbf{Research Intern}  \hfill {\em May 2020 - August 2020} \\
Microsoft Research, Redmond, US \\
Internship in the RiSE group; advised by Sebastian Burckhardt. \\
Worked on an execution engine for stateful serverless applications.

\textbf{Software Engineering Intern}  \hfill {\em May 2019 - August 2019} \\
Amazon Web Services, New York, US \\
Internship in the Automated Reasoning Group; advised by Daniel Schwartz-Narbonne. \\
Worked on the verification of critical C code.

\textbf{Google Summer of Code Student}  \hfill {\em May 2017 - September 2017} \\
Worked with BEAM Community to extend the \textit{ejabberd} open source project. \\
Implemented support for the \textit{``Let's Encrypt'' ACME} certificate acquiring protocol.

\textbf{Big Data Application Developer}  \hfill {\em June 2016 - September 2016} \\
Everis, Barcelona, Spain \\
Internship at the Big Data Center of Excellence. \\
Developed Big Data Applications using a plethora of tools in the Hadoop ecosystem.




\section{Teaching Experience}

\textbf{Teaching Assistant}  \hfill {\em Fall 2021} \\
Institution: University of Pennsylvania \\
Course: \textit{Computer-Aided Verification}, Graduate level \\
Professor: Rajeev Alur

\textbf{Teaching Assistant}  \hfill {\em Fall 2019} \\
Institution: University of Pennsylvania \\
Course: \textit{Software Foundations}, Graduate level \\
Professor: Benjamin Pierce

\textbf{Lab Assistant}  \hfill {\em Fall 2017} \\
Institution: National Technical University of Athens \\
Course: \textit{Introduction to Programming}, Undergraduate level \\
Professors: S. Zachos, N. Papaspyrou, V. Kantere, and P. Potikas

\section{Publications}

\begin{minipage}{\textwidth}
\textbf{Practically Correct, Just-in-Time Shell Script Parallelization}. \\
Konstantinos Kallas, Tammam Mustafa, Jan Bielak, Dimitris Karnikis, Thurston H.Y. Dang, Michael Greenberg, and Nikos Vasilakis. \\
16th USENIX Symposium on Operating Systems Design and Implementation (OSDI 22).
\end{minipage}

\begin{minipage}{\textwidth}
\textbf{Netherite: Efficient Execution of Serverless Workflows}. \\
Sebastian Burckhardt, Badrish Chandramouli, Chris Gillum, David Justo, Konstantinos Kallas, Connor McMahon, Christopher S. Meiklejohn, and Xiangfeng Zhu. \\
Proceedings of the VLDB Endowment (VLDB 2022).
\end{minipage}

\begin{minipage}{\textwidth}
\textbf{Stream Processing with Dependency-Guided Synchronization}. \\
Konstantinos Kallas$^*$, Filip Niksic$^*$, Caleb Stanford$^*$, and Rajeev Alur. \\
Proceedings of the 27th ACM SIGPLAN Symposium on Principles and Practice of Parallel Programming (PPoPP 2022).
\end{minipage}

\begin{minipage}{\textwidth}
\textbf{Charon: A Framework for Microservice Overload Control}. \\
Jiali Xing, Max Demoulin, Konstantinos Kallas, and Benjamin C. Lee. \\
Proceedings of the 18th ACM Workshop on Hot Topics in Networks (HotNets 2021).
\end{minipage}

\begin{minipage}{\textwidth}
\textbf{Durable Functions: Semantics for Stateful Serverless}. \\
Sebastian Burckhardt, Chris Gillum, David Justo, Konstantinos Kallas, Connor McMahon, and Christopher S. Meiklejohn. \\
Proceedings of the ACM on Programming Languages (OOPSLA 2021).
\end{minipage}

\begin{minipage}{\textwidth}
\textbf{An Order-aware Dataflow Model for Parallel Unix Pipelines}. \\
Shivam Handa$^*$, Konstantinos Kallas$^*$, Nikos Vasilakis$^*$, and Martin Rinard. \\
Proceedings of the ACM on Programming Languages (ICFP 2021).
\end{minipage}

\begin{minipage}{\textwidth}
\textbf{Synchronization Schemas}. \\
Rajeev Alur, Phillip Hillard, Zachary G. Ives, Konstantinos Kallas, Konstantinos Mamouras, Filip Niksic, Caleb Stanford, Val Tannen, and Anton Xue. \\
Invited Paper at Proceedings of the 40th Symposium on Principles of Database Systems (PODS 2021).
\end{minipage}

\begin{minipage}{\textwidth}
\textbf{Unix Shell Programming: The Next 50 Years}. \\
Michael Greenberg$^*$, Konstantinos Kallas$^*$, and Nikos Vasilakis$^*$. \\
Proceedings of the Workshop on Hot Topics in Operating Systems (HotOS 2021).\\
 \emph{Distinguished Presentation Award.}
\end{minipage}

\begin{minipage}{\textwidth}
\textbf{The Future of the Shell: Unix and Beyond}. \\
Michael Greenberg$^*$, Konstantinos Kallas$^*$, and Nikos Vasilakis$^*$. \\
Panel at the Workshop on Hot Topics in Operating Systems (HotOS 2021).
\end{minipage}

\begin{minipage}{\textwidth}
\textbf{PaSh: Light-touch Data-Parallel Shell Processing}. \\
Nikos Vasilakis$^*$, Konstantinos Kallas$^*$, Konstantinos Mamouras, Achilleas Benetopoulos, and Lazar M. Cvetković. \\
Proceedings of the Sixteenth European Conference on Computer Systems (EuroSys 2021).\\
 \emph{Best Paper Award.}
\end{minipage}

\begin{minipage}{\textwidth}
\textbf{Preventing Dynamic Library Compromise on Node. js via RWX-Based Privilege Reduction}. \\
Nikos Vasilakis, Cristian-Alexandru Staicu, Grigoris Ntousakis, Konstantinos Kallas, Ben Karel, André DeHon, and Michael Pradel. \\
Proceedings of the ACM SIGSAC Conference on Computer and Communications Security (CCS’21).
\end{minipage}

\begin{minipage}{\textwidth}
\textbf{Code-level model checking in the software development workflow at Amazon Web Services}. \\
Nathan Chong, Byron Cook, Jonathan Eidelman, Konstantinos Kallas, Kareem Khazem, Felipe R. Monteiro, Daniel Schwartz-Narbonne, Serdar Tasiran, Michael Tautschnig, and Mark R. Tuttle. \\
Software: Practice and Experience 2021.
\end{minipage}

\begin{minipage}{\textwidth}
\textbf{DiffStream: Differential Output Testing for Stream Processing Programs}. \\
Konstantinos Kallas$^*$, Filip Niksic$^*$, Caleb Stanford$^*$, and Rajeev Alur. \\
Proceedings of the ACM on Programming Languages (OOPSLA 2020).
\end{minipage}

\begin{minipage}{\textwidth}
\textbf{Code-Level Model Checking in the Software Development Workflow}. \\
Nathan Chong, Byron Cook, Konstantinos Kallas, Kareem Khazem, Felipe R. Monteiro, Daniel Schwartz-Narbonne, Serdar Tasiran, Michael Tautschnig, and Mark R. Tuttle. \\
42st International Conference on Software Engineering: Software Engineering in Practice (ICSE-SEIP 2020).
\end{minipage}

\begin{minipage}{\textwidth}
\textbf{Security Criteria for a Transparent Encryption Layer}. \\
Konstantinos Kallas, Clara Schneidewind, Benjamin C. Pierce, and Steve Zdancewic. \\
Workshop on Foundations of Computer Security (FCS 2019).
\end{minipage}

\begin{minipage}{\textwidth}
\textbf{HiPErJiT: A Profile-Driven Just-in-Time Compiler for Erlang}. \\
Konstantinos Kallas and Konstantinos Sagonas. \\
30th Symposium on Implementation and Application of Functional Languages (IFL 2018).
\end{minipage}




Note: $^*$ indicates equal contribution.

\section{Honors and Awards}

\textbf{A.G. Leventis Foundation PhD Grant} \hfill {\em 2021-2023}

\textbf{ACM SRC Grand Finals} \hfill {\em 2021} \\
\textit{2nd place} among SRC winners across all ACM conferences.

\textbf{HotOS 2021 Distinguished Presentation Award} \hfill {\em 2021} \\
Awarded for "Unix Shell Programming: The Next 50 Years".

\textbf{EuroSys 2021 Best Paper Award} \hfill {\em 2021} \\
Awarded for "PaSh: Light-touch Data-Parallel Shell Processing".

\textbf{POPL Student Research Competition} \hfill {\em 2021} \\
\textit{1st place} at the graduate category of the research competition. \\
Presented work on a parallelizing JiT compiler for shell scripts. 

\textbf{Gerondelis Foundation PhD Award} \hfill {\em 2020}

\textbf{Programming Competitions}  \hfill {\em 2015-2018} \\
Participation in many programming competitions. Notable examples: \\
ICFP Programming Contest 2018 (Lightning) \textit{10th out of 91 teams}, \\
IEEE Xtreme 2017 \textit{Top 5\%} and \textit{3rd} in Greece, IEEE Xtreme 2016 \textit{Top 10\%}

\textbf{Heterogenous Computing Student Challenge Certificate}  \hfill {\em 2017} \\
HiPEAC CSW, Zagreb, Croatia \\
Optimizing GPU implementation of the K-means algorithm (NTUA-team) \\
Supervision: Prof. Georgios Goumas

\textbf{EESTech Challenge}  \hfill {\em 2017} \\ 
Supervised Machine Learning Hackathon \\
\textit{Joint 1st place} among 40 teams.

\textbf{The Great Moment of Education Scholarship from Eurobank EFG}  \hfill {\em 2012} \\ 
Achieving the highest rank in national qualifications exams in Dionisos high school.

\textbf{9th European Union Science Olympiad} \hfill {\em 2011} \\ 
Team-based science competition in Biology, Chemistry, and Physics.\\
\textit{1st place} in local round and \textit{3rd place} in  national round. 

\textbf{Mathematical Competition, Hellenic Mathematical Society} \hfill {\em 2010} \\ 
Mathematical competition for high school students. \\
Distinction in the 1st and 2nd local round. 


%% \section{Selected University Projects}

%% \textbf{Edsger Compiler} \href{https://github.com/angelhof/edsger_compiler}{\includegraphics[width=0.31cm,height=0.31cm,keepaspectratio,trim=0 0.13cm 0 0]{GitHub-Mark-32px}}\hfill {\em Spring 2016} \\
%% Course: \textit{Compilers} -- Professors: K. Sagonas and N. Papaspyrou \\
%% Complete compiler for Edsger (a C-like educational language). \\
%% Implemented a lexer, a parser, a semantic analysis pass and a transpiler to LLVM-IR.

%% \textbf{Chord Protocol Implementation} \href{https://github.com/angelhof/chord-distributed-erlang}{\includegraphics[width=0.31cm,height=0.31cm,keepaspectratio,trim=0 0.13cm 0 0]{GitHub-Mark-32px}} \hfill {\em Fall 2016} \\
%% Course: \textit{Distributed Systems} -- Professor: N. Koziris \\
%% Implemented two versions of the Chord protocol, one with sequential replica consistency and one with eventual replica consistency. \\
%% Tested whether the replicated values satisfied the constraints set by the consistency protocols.

%% \textbf{Song Emotion Classifier} \href{https://github.com/etzinis/pat_rec_ntua}{\includegraphics[width=0.31cm,height=0.31cm,keepaspectratio,trim=0 0.13cm 0 0]{GitHub-Mark-32px}} \hfill {\em Fall 2016} \\
%% Course: \textit{Pattern Recognition} -- Professor: A. Potamianos \\
%% Worked on an automatic emotion recognition system, dedicated to classifying the emotional content of songs. \\
%% The system extracted the essential features from the songs and trained an SVM classifier.


% \section{Computer Experience}

% \textbf{Programming Languages:} \\
% Erlang, Haskell, Python, Coq, C, Scala, SML, Agda, Javascript, Matlab, Java, SQL
% \\
% \textbf{Operating Systems:} \\
% Extensive experience on the Linux OS and Linux Kernel Programming.
% \\
% \textbf{Hadoop:} \\
% Deep understanding of the Hadoop Ecosystem as well as hands on experience.
% \\
% \textbf{Typesetting:} \\
% Experience using \LaTeX.  %% \space for writing reports and presentations 



\section{Languages}

\textbf{Greek} (Native), \textbf{English} (C2), \textbf{Italian} (C1), \textbf{German} (B2)


%% \section{Interests}

%% \textbf{Pastimes:} \\
%% Guitar playing and music in general, Board Games, Table Tennis, Cinema, Travelling \\
%% \textbf{Programming:} \\
%% Working on several hobby and open source projects. \\
%% \textbf{Memberships:} \\
%% \textit{IEEE} Member, \textit{ESN} Member \\

\end{resume}
\end{document}
